Il presente lavoro ha affrontato il problema della pianificazione della raccolta differenziata urbana attraverso un modello di ottimizzazione bi-obiettivo, in cui la qualità del servizio percepita dai cittadini e l'efficienza economica dell'azienda di gestione vengono ricondotte a un'unica funzione obiettivo $F_{tot}$, bilanciata dal parametro decisionale $\gamma$. Il modello è stato risolto attraverso due euristiche di instradamento concettualmente distinte, un approccio costruttivo competitivo multi-veicolo (Greedy) e un approccio basato sulla matrice dei risparmi (Clarke-Wright), entrambe integrate con una procedura di ricerca a griglia per la selezione della frequenza di raccolta ottima $X_r^*$ per ciascuna tipologia di rifiuto.

La campagna sperimentale, articolata in cinque benchmark specifici e due benchmark trasversali, ha permesso di isolare progressivamente l'effetto delle singole variabili del problema, per poi ricomporle in un quadro interpretativo unitario. Dall'insieme dei risultati emergono tre evidenze principali.

In primo luogo, il parametro $\gamma$ si conferma il vero regolatore strutturale del sistema. La sua variazione non si limita a spostare il valore numerico di $F_{tot}$, ma trasforma qualitativamente la natura del vincolo operativo attivo: nelle politiche orientate all'efficienza economica il collo di bottiglia è la capacità volumetrica della flotta, mentre nelle politiche orientate alla qualità del servizio il vincolo diventa la durata massima del turno lavorativo. La politica Neutra ($\gamma = 0.50$) si colloca costantemente, in tutti i benchmark analizzati, in una posizione di compromesso favorevole, in cui gran parte dell'insoddisfazione degli utenti viene assorbita a fronte di un incremento dei costi logistici ancora proporzionalmente sostenibile.

In secondo luogo, il confronto sistematico tra le due euristiche ha evidenziato una superiorità strutturale e non occasionale dell'algoritmo Clarke-Wright. La dominanza in senso paretiano osservata nel Benchmark Trasversale 1, unita alla sistematicità del vantaggio riscontrato su istanze sintetiche di diversa dimensione, tipologia urbanistica e distribuzione spaziale, nonché sulla mappa reale di Fabriano, indica che tale vantaggio non dipende dalle caratteristiche specifiche di una singola istanza, ma da una reale differenza qualitativa nella strategia di costruzione delle rotte: la valutazione globale delle opportunità di fusione consente alla Clarke-Wright di individuare configurazioni di routing più compatte di quanto la logica sequenziale e locale del Greedy sia in grado di ottenere, pur rimanendo quest'ultimo un'euristica competitiva grazie al meccanismo di bilanciamento del carico introdotto dalla valutazione della funzione obiettivo parziale a ogni passo.

In terzo luogo, la formulazione disaccoppiata delle frequenze di raccolta per tipologia di rifiuto si è rivelata essenziale per catturare la reale eterogeneità del problema. L'organico e l'indifferenziato, gravati da un'elevata sensibilità dell'utenza all'accumulo, richiedono frequenze di ritiro sensibilmente superiori rispetto a frazioni a bassa putrescibilità come il vetro; analogamente, l'analisi di sensibilità condotta nel Benchmark Trasversale 2 ha mostrato come il legame tra frequenza di raccolta e dimensionamento della flotta sia tutt'altro che intuitivo, con le configurazioni a bassissima frequenza che, anziché ridurre il numero di veicoli necessari, lo massimizzano a causa della saturazione fisica generata dall'accumulo dei rifiuti presso le utenze.

Nel complesso, il sistema sviluppato si configura non come un semplice risolutore che produce un'unica soluzione ottima, ma come un vero e proprio strumento di supporto alle decisioni, capace di esplicitare in modo quantitativo il costo economico di un determinato livello di servizio e, specularmente, l'impatto sociale di una politica orientata al contenimento della spesa.

\section*{Limiti del modello}
\addcontentsline{toc}{section}{Limiti del modello}

Accanto ai risultati ottenuti, è opportuno individuare con onestà i limiti della formulazione proposta, alcuni intrinseci alle scelte di modellazione, altri legati alle euristiche di risoluzione adottate.

Un primo limite riguarda la natura dei dati di domanda. La produzione di rifiuto $W$ e la frequenza ideale $x^*_{r,u}$ associate a ciascuna tipologia di utenza sono state definite come parametri costanti e deterministici, stimati a priori e non calibrati su serie storiche di raccolta reali. Tale semplificazione vale anche per il benchmark condotto sulla mappa di Fabriano, nel quale la componente realistica dell'istanza riguarda esclusivamente la struttura della rete stradale e la distribuzione geografica e tipologica degli utenti, mentre le quantità di rifiuto prodotte continuano a derivare dagli stessi parametri sintetici utilizzati nelle istanze artificiali. Il modello non incorpora inoltre alcuna forma di variabilità stocastica o stagionale della domanda, né un margine di incertezza sulle stime di produzione.

Un secondo limite riguarda l'assenza di un riferimento esatto rispetto al quale misurare la qualità assoluta delle soluzioni. Il confronto condotto in questo lavoro è sempre relativo, tra le due euristiche implementate, senza il supporto di un lower bound analitico o di una soluzione esatta ottenuta tramite programmazione lineare intera o algoritmi di branch-and-bound su istanze di dimensione contenuta. Il gap prestazionale quantificato nei diversi benchmark descrive quindi la distanza tra Greedy e Clarke-Wright, non la distanza di entrambi dall'ottimo globale del problema.

Un terzo limite è di natura strutturale e riguarda il disaccoppiamento completo tra le cinque tipologie di rifiuto. Ciascun flusso di raccolta viene ottimizzato in modo indipendente, con una flotta dedicata e un instradamento proprio, senza considerare la possibilità di una raccolta congiunta mediante veicoli a scomparti multipli, pratica sempre più diffusa nella gestione reale dei servizi di igiene urbana. Analogamente, il modello assume una flotta omogenea per ciascuna tipologia di rifiuto, senza la possibilità di combinare veicoli di capacità differente all'interno della stessa categoria.

Infine, il fattore di scala $k_{scala}$, necessario per rendere confrontabili la componente di insoddisfazione e la componente logistica all'interno della funzione obiettivo, viene stimato attraverso un'unica esecuzione preliminare dell'euristica Greedy sull'intera griglia di frequenze. Questa scelta implementativa, per quanto giustificata da esigenze di efficienza computazionale, introduce una dipendenza tra la normalizzazione dell'obiettivo e le caratteristiche specifiche dell'algoritmo utilizzato per calcolarla.

\section*{Possibili sviluppi futuri}
\addcontentsline{toc}{section}{Possibili sviluppi futuri}

I limiti individuati suggeriscono direttrici di sviluppo naturali per proseguire ed estendere il lavoro svolto, distinguibili in tre ambiti principali.

\subsection*{Metaeuristiche avanzate}
\addcontentsline{toc}{subsection}{Metaeuristiche avanzate}

Le due euristiche costruttive impiegate in questo lavoro, pur avendo dimostrato buone proprietà di scalabilità grazie alla vettorizzazione numerica e alle strutture dati a complessità O(1), non includono alcuna fase di miglioramento locale successiva alla costruzione iniziale della soluzione. Un naturale sviluppo consiste nell'applicare operatori di ricerca locale, quali lo scambio 2-opt o l'or-opt, come fase di post-elaborazione delle rotte prodotte da Greedy e Clarke-Wright, oppure nell'adottare metaeuristiche più strutturate, quali il Simulated Annealing, la Tabu Search o l'Adaptive Large Neighborhood Search, capaci di sfuggire agli ottimi locali tipici delle euristiche costruttive. Parallelamente, l'attuale esplorazione del compromesso tra le due componenti dell'obiettivo, limitata a tre soli valori discreti di $\gamma$, potrebbe essere estesa mediante algoritmi metaeuristici multi-obiettivo, come NSGA-II, in grado di ricostruire l'intera frontiera di Pareto in modo continuo anziché per punti isolati, fornendo un supporto decisionale ancora più granulare ai pianificatori del servizio.

\subsection*{Dati dinamici e traffico reale}
\addcontentsline{toc}{subsection}{Dati dinamici e traffico reale}

Il modello attuale assume matrici di distanza e di tempo statiche, calcolate a partire dalla sola geometria della rete stradale e da una velocità media costante, senza tenere conto della variabilità del traffico in funzione della fascia oraria o del giorno della settimana. Un'estensione rilevante consiste nell'integrazione di dati di traffico in tempo reale, ad esempio tramite API di provider cartografici, per ottenere matrici di tempo time-dependent e abilitare strategie di ri-pianificazione dinamica durante il turno di lavoro. Analogamente, la sostituzione dei parametri di produzione del rifiuto $W$, attualmente stimati in modo statico, con letture provenienti da sensori di riempimento installati sui cassonetti (approccio smart-bin, tipico dei sistemi IoT per la gestione dei rifiuti) permetterebbe di alimentare il modello con dati di domanda realmente osservati, superando l'ipotesi di produzione costante e introducendo la possibilità di modellare esplicitamente la stagionalità dei conferimenti.

\subsection*{Integrazione con sistemi intelligenti}
\addcontentsline{toc}{subsection}{Integrazione con sistemi intelligenti}

Da ultimo, il sistema sviluppato si presta a essere integrato all'interno di piattaforme più ampie di gestione intelligente del territorio. L'adozione di modelli di apprendimento automatico per la previsione della domanda di raccolta a livello di singola utenza, a partire da serie storiche e da variabili esogene quali la composizione demografica o la stagionalità, potrebbe raffinare ulteriormente la stima dei parametri $W$ e $x^*_{r,u}$ oggi assunti come costanti. In una prospettiva più ambiziosa, tecniche di reinforcement learning potrebbero essere impiegate per addestrare agenti di instradamento capaci di adattare le proprie decisioni in modo continuo alle condizioni operative osservate, superando la rigidità di una pianificazione basata su euristiche statiche. Infine, l'integrazione dei risultati del modello all'interno di una dashboard decisionale, o di un vero e proprio gemello digitale (digital twin) del sistema di raccolta, consentirebbe agli amministratori pubblici e ai gestori del servizio di esplorare interattivamente lo spazio delle soluzioni, valutando in tempo reale l'impatto di scelte alternative sulla qualità del servizio e sui costi operativi complessivi.